% =============================================================================
% Manual P01 – Protección Ocular frente a Radiación Láser
% Técnicas de Medida – Ingeniería Aeronáutica
% Compilar con: latexmk -xelatex -interaction=nonstopmode manual_P01.tex
% =============================================================================
\documentclass[12pt,a4paper]{article}

% ---- Encoding & Language ----
\usepackage{fontspec}
\usepackage[spanish]{babel}
\addto\captionsspanish{\renewcommand{\tablename}{Tabla}}

% ---- Page Layout ----
\usepackage[top=2.5cm,bottom=2.5cm,left=2.5cm,right=2.5cm]{geometry}
\setlength{\headheight}{14.5pt}

% ---- Math & Physics ----
\usepackage{amsmath,amssymb,amsfonts}
\usepackage{mathtools}
\usepackage{siunitx}
\sisetup{output-decimal-marker={,}, group-separator={.}}

% ---- Graphics & Color ----
\usepackage{graphicx}
\usepackage[table]{xcolor}
\usepackage{float}
\usepackage{caption}
\usepackage{subcaption}

% ---- TikZ ----
\usepackage{tikz}
\usetikzlibrary{babel}
\usetikzlibrary{arrows.meta,positioning,calc,shapes.geometric,decorations.pathmorphing,
                decorations.markings,patterns,fit,backgrounds}

% ---- Tables ----
\usepackage{booktabs}
\usepackage{multirow}
\usepackage{array}
\usepackage{colortbl}
\usepackage{tabularx}
\usepackage{adjustbox}
\usepackage{longtable}

% ---- Code Listings ----
\usepackage{listings}
\lstset{
  basicstyle=\ttfamily\small,
  keywordstyle=\color{blue!70!black}\bfseries,
  commentstyle=\color{green!50!black},
  stringstyle=\color{red!60!black},
  backgroundcolor=\color{gray!10},
  frame=single,
  breaklines=true,
  numbers=left,
  numberstyle=\tiny\color{gray}
}

% ---- Hyperlinks ----
\usepackage[colorlinks=true,linkcolor=blue!70!black,citecolor=green!50!black,
            urlcolor=blue!50!black]{hyperref}

% ---- Bibliography ----
\usepackage[backend=biber,style=ieee,sorting=none]{biblatex}
\addbibresource{references_P01.bib}

% ---- Headers & Footers ----
\usepackage{fancyhdr}
\pagestyle{fancy}
\fancyhf{}
\fancyhead[L]{\small Técnicas de Medida -- Ingeniería Aeronáutica}
\fancyhead[R]{\small P01: Protección Láser}
\fancyfoot[C]{\thepage}

% ---- Custom Colors ----
\definecolor{aeroblue}{RGB}{0,51,102}
\definecolor{aerored}{RGB}{153,0,0}
\definecolor{aerogreen}{RGB}{0,102,51}
\definecolor{aeroyellow}{RGB}{255,204,0}

% =============================================================================
\begin{document}

\begin{titlepage}
	\centering
	\vspace*{1cm}
	{\color{aeroblue}\rule{\linewidth}{2pt}}\\[0.4cm]
	{\Large\bfseries\color{aeroblue} TÉCNICAS DE MEDIDA}\\[0.3cm]
	{\color{aeroblue}\rule{\linewidth}{1pt}}\\[1.5cm]
	{\huge\bfseries P01 -- Protección Ocular\\[0.3cm]
	frente a Radiación Láser}\\[1cm]
	{\large Manual de Laboratorio}\\[0.5cm]
	{\large\itshape Ingeniería Aeronáutica}\\[2cm]
	\begin{tikzpicture}
		\draw[fill=aeroblue!20,draw=aeroblue,thick] (0,0) rectangle (1.5,0.6);
		\node at (0.75,0.3) {\small\bfseries LÁSER};
		\draw[aerored,ultra thick,->] (1.5,0.3) -- (5.5,0.3);
		\draw[aerored!40,thin] (1.5,0.3) -- (5.5,0.8);
		\draw[aerored!40,thin] (1.5,0.3) -- (5.5,-0.2);
		\draw[fill=gray!30,draw=black] (5.5,-0.15) ellipse (0.35 and 0.55);
		\draw[aerogreen,very thick] (5.1,0) rectangle (5.9,0.6);
		\node[aerogreen,font=\tiny\bfseries] at (5.5,0.3) {OD};
		\node[below,font=\footnotesize,aerored] at (3.5,-0.1) {$\lambda,\; P,\; E$};
	\end{tikzpicture}\\[2.5cm]
	{\color{aeroblue}\rule{\linewidth}{1pt}}
\end{titlepage}

\tableofcontents
\newpage

% =============================================================================
\section{Introducción}
\label{sec:intro}

Los láseres (Amplificación de Luz por Emisión Estimulada de Radiación) son
herramientas fundamentales en numerosos campos de la ingeniería aeronáutica,
especialmente en laboratorios de investigación y desarrollo para técnicas de
medición avanzadas como la Anemometría Láser Doppler (LDA), la Velocimetría
por Imagen de Partículas (PIV), la metrología de precisión y los ensayos no
destructivos~\cite{svelto2010}.

La naturaleza altamente concentrada y coherente de la radiación láser, si bien
es ventajosa para estas aplicaciones, representa un riesgo significativo para la
seguridad, particularmente para la visión~\cite{henderson2004}.

Esta práctica introduce los principios fundamentales de la seguridad láser, con
un enfoque en la comprensión analítica de los riesgos asociados a diferentes
tipos de láseres comúnmente empleados en el contexto aeronáutico. Se abordará
la selección y especificación del Equipo de Protección Ocular (EPO) adecuado,
a través del análisis de datos técnicos, el cálculo de parámetros de seguridad
críticos y la toma de decisiones informada.

% =============================================================================
\section{Objetivos de Aprendizaje}
\label{sec:objetivos}

Al finalizar esta práctica, el estudiante será capaz de:
\begin{enumerate}
	\item Identificar y describir los peligros asociados a la radiación láser en
	      técnicas de medida y aplicaciones aeronáuticas.
	\item Comprender y aplicar la clasificación de seguridad láser según
	      ANSI~Z136.1~\cite{ansi2014} e IEC~60825-1~\cite{iec60825}.
	\item Analizar cualitativamente la interacción de la radiación láser con el
	      tejido ocular en función de la longitud de onda.
	\item Calcular la Exposición Potencial ($H_0$) a partir de los parámetros
	      de una fuente láser.
	\item Calcular la Densidad Óptica (OD) mínima requerida, dados los
	      parámetros del láser y la Exposición Máxima Permisible (MPE).
	\item Seleccionar el equipo de protección ocular apropiado de un catálogo
	      de 12 lentes, justificando la elección mediante cálculos.
	\item Describir los elementos fundamentales de un programa de seguridad
	      láser en un entorno de laboratorio aeronáutico.
\end{enumerate}

% =============================================================================
\section{Fundamentos Teóricos}
\label{sec:teoria}

\subsection{Principios de la Radiación Láser}

La radiación láser se distingue de la luz convencional por cuatro propiedades
fundamentales~\cite{svelto2010,saleh2007}:

\begin{enumerate}
	\item \textbf{Monocromaticidad:} luz de una única longitud de onda.
	\item \textbf{Coherencia:} ondas en fase espacial y temporalmente.
	\item \textbf{Direccionalidad (Colimación):} baja divergencia.
	\item \textbf{Alta Irradiancia:} gran concentración energética.
\end{enumerate}

\subsection{Clasificación de Seguridad Láser}

Los estándares ANSI~Z136.1~\cite{ansi2014} e IEC~60825-1~\cite{iec60825}
clasifican los láseres según su potencial de causar daño
(Tabla~\ref{tab:clases}).

\begin{table}[H]
	\centering
	\caption{Clasificación de seguridad láser (IEC~60825-1 / ANSI~Z136.1).}
	\label{tab:clases}
	\small
	\begin{tabularx}{\textwidth}{c l l X}
		\toprule
		\textbf{Clase} & \textbf{Riesgo}                              & \textbf{Ejemplo}           & \textbf{Controles}        \\
		\midrule
		1              & Seguro en uso normal                         & Sensores encapsulados      & Ninguno especial          \\
		1M             & Seguro a simple vista                        & Fibra óptica con lente     & Sin óptica de aumento     \\
		2              & Visible, reflejo palpebral ($P<\SI{1}{\mW}$) & Puntero                    & Señalización              \\
		2M             & Como Clase~2, sin óptica de aumento          & {--}                       & Señalización              \\
		3R             & Riesgo bajo directa ($P<\SI{5}{\mW}$)        & Diodo metrología           & Señalización + precaución \\
		3B             & Peligroso directa ($P<\SI{500}{\mW}$)        & LDA 50~mW                  & EPO frecuente             \\
		\rowcolor{aerored!20}
		\textbf{4}     & \textbf{Directa, especular y difusa}         & \textbf{PIV, LDA alta $P$} & \textbf{EPO obligatorio}  \\
		\bottomrule
	\end{tabularx}
\end{table}

\subsection{Interacción Láser–Tejido Ocular}

El daño ocular depende críticamente de la longitud de onda~\cite{henderson2004}:

\begin{itemize}
	\item \textbf{UV (100--400~nm):} córnea/cristalino $\rightarrow$ fotoqueratitis.
	\item \textbf{Visible + IR-A (400--1400~nm):} retina con factor $\sim\!10^5$
	      $\rightarrow$ \textbf{zona de máximo peligro retiniano}.
	\item \textbf{IR-B/C ($>$1400~nm):} córnea $\rightarrow$ quemaduras.
\end{itemize}

\begin{figure}[H]
	\centering
	\begin{tikzpicture}[scale=0.9]
		\fill[violet!40] (0,0) rectangle (2,1);
		\fill[yellow!40] (2,0) rectangle (5,1);
		\fill[red!20]    (5,0) rectangle (8,1);
		\node[font=\footnotesize,rotate=90] at (1,0.5) {UV};
		\node[font=\footnotesize,rotate=90] at (3.5,0.5) {Visible + IR-A};
		\node[font=\footnotesize,rotate=90] at (6.5,0.5) {IR-B / IR-C};
		\draw[aerored,very thick] (0,1.5) -- (1,1.2) -- (2,2.5) -- (3.5,3.5)
		-- (5,2.5) -- (6,1.5) -- (8,1.2);
		\node[aerored,font=\footnotesize,above] at (3.5,3.5) {Máximo peligro retiniano};
		\draw[->] (0,0) -- (8.5,0) node[right,font=\footnotesize] {$\lambda$~(nm)};
		\node[font=\tiny] at (0,-0.3) {100};
		\node[font=\tiny] at (2,-0.3) {400};
		\node[font=\tiny] at (5,-0.3) {1400};
		\node[font=\tiny] at (8,-0.3) {$10^6$};
		\draw[aerored,dashed,thick] (3.2,0) -- (3.2,3.0);
		\fill[aerored] (3.2,3.0) circle (0.08);
		\node[aerored,font=\tiny,right] at (3.2,3.0) {532~nm (PIV)};
		\draw[blue!70,dashed,thick] (3.0,0) -- (3.0,2.8);
		\fill[blue!70] (3.0,2.8) circle (0.08);
		\node[blue!70,font=\tiny,left] at (3.0,2.8) {514,5~nm (LDA)};
		\draw[red!70,dashed,thick] (3.8,0) -- (3.8,3.3);
		\fill[red!70] (3.8,3.3) circle (0.08);
		\node[red!70,font=\tiny,right] at (3.8,3.2) {635~nm (Alin.)};
	\end{tikzpicture}
	\caption{Absorción espectral en el ojo humano y ubicación de los tres láseres
		de esta práctica.}
	\label{fig:espectro}
\end{figure}

\subsection{Exposición Máxima Permisible (MPE)}

La MPE es el nivel de radiación láser al cual una persona puede estar expuesta
sin sufrir efectos adversos~\cite{iec60825,ansi2014}. Depende de $\lambda$,
duración de exposición y tejido expuesto.

\subsection{Densidad Óptica (OD)}
\label{sec:OD_teoria}

\begin{equation}
	\boxed{OD = \log_{10}\!\left(\frac{H_0}{H_{\mathrm{MPE}}}\right)}
	\label{eq:OD}
\end{equation}
donde $H_0$ es la exposición potencial sin protección y $H_{\mathrm{MPE}}$ la
MPE. Una OD de $N$ atenúa por factor $10^N$.

\subsection{Transmisión de Luz Visible (VLT)}

VLT (\%) = porcentaje de luz visible (400--700~nm) transmitida. Mayor VLT =
mejor visibilidad en laboratorio. Existe compromiso OD↑ $\Leftrightarrow$
VLT↓.

% =============================================================================
\section{Escenarios Láser y Datos de la Práctica}
\label{sec:escenarios}

\subsection{Fuentes Láser}

Se analizan tres escenarios representativos (Tabla~\ref{tab:lasers}).

\begin{table}[H]
	\centering
	\caption{Parámetros de las fuentes láser de la práctica.}
	\label{tab:lasers}
	\small
	\begin{tabular}{l c c c}
		\toprule
		\textbf{Parámetro} & \textbf{L1 -- PIV} & \textbf{L2 -- LDA} & \textbf{L3 -- Alin.} \\
		\midrule
		Tipo               & Nd:YAG $2\omega$   & Argón Ion CW       & Diodo CW             \\
		$\lambda$ (nm)     & 532                & 514,5              & 635                  \\
		Modo               & Pulsado            & CW                 & CW                   \\
		$E$ (mJ/pulso)     & 200                & {--}               & {--}                 \\
		$P$ (W)            & {--}               & 1,5                & 0,0045               \\
		$\tau$ (ns)        & 8                  & {--}               & {--}                 \\
		$f$ (Hz)           & 15                 & {--}               & {--}                 \\
		$a$ (mm)           & 5                  & 1,2                & 3                    \\
		\rowcolor{aerored!15}
		\textbf{Clase}     & \textbf{4}         & \textbf{4}         & \textbf{3R}          \\
		\bottomrule
	\end{tabular}
\end{table}

\subsection{Catálogo de Gafas EPO (12 lentes)}
\label{sec:catalogo}

El catálogo ampliado incluye 12 modelos de gafas de protección ocular
(Tabla~\ref{tab:epo}). Para cada lente se indica la banda de protección
(rango $\lambda$), la OD en esa banda y la transmisión de luz visible.

\begin{table}[H]
	\centering
	\caption{Catálogo completo de gafas de protección ocular (12 lentes).}
	\label{tab:epo}
	\small
	\begin{tabularx}{\textwidth}{c l l X c}
		\toprule
		\textbf{ID} & \textbf{Marca/Modelo}                                & \textbf{Lente}
		            & \textbf{Bandas de protección (OD @ $\lambda$)}
		            & \textbf{VLT}                                                           \\
		\midrule
		L-01        & Thorlabs LG1                                         & Naranja
		            & OD~5+ @ 190--540~nm
		            & 40\%                                                                   \\
		L-02        & Thorlabs LG7                                         & Verde
		            & OD~2+ @ 630--670~nm; OD~4+ @ 800--1100~nm
		            & 55\%                                                                   \\
		L-03        & Laservision F18                                      & Azul cobalto
		            & OD~7+ @ 190--400~nm; OD~6+ @ 532~nm; OD~5+ @ 1064~nm
		            & 15\%                                                                   \\
		L-04        & Kentek KXL-62C                                       & Transparente
		            & OD~2+ @ 10\,600~nm (CO$_2$)
		            & 90\%                                                                   \\
		L-05        & Laservision P1H03                                    & Multi-$\lambda$
		            & OD~7 @ 190--534~nm; OD~6 @ 800--820~nm
		            & 30\%                                                                   \\
		L-06        & Thorlabs LG3                                         & Ámbar
		            & OD~3+ @ 190--532~nm
		            & 50\%                                                                   \\
		L-07        & NoIR LaserShields                                    & Rojo oscuro
		            & OD~4+ @ 488--540~nm
		            & 35\%                                                                   \\
		L-08        & Honeywell GPT                                        & Gris
		            & OD~3+ @ 630--694~nm; OD~5+ @ 800--1100~nm
		            & 25\%                                                                   \\
		L-09        & Uvex 9116                                            & Naranja claro
		            & OD~2+ @ 315--532~nm; OD~4+ @ 808--1100~nm
		            & 45\%                                                                   \\
		L-10        & Kentek KXL-55A                                       & Verde jade
		            & OD~6+ @ 190--540~nm
		            & 20\%                                                                   \\
		L-11        & NoIR DI2                                             & Rosa
		            & OD~1+ @ 585--650~nm
		            & 60\%                                                                   \\
		L-12        & Laservision R14T1C04                                 & Multi-$\lambda$
		            & OD~4+ @ 190--534~nm; OD~5+ @ 630--700~nm
		            & 22\%                                                                   \\
		\bottomrule
	\end{tabularx}
\end{table}

\textbf{Nota:} los archivos \texttt{data/epo\_lenses.csv} y
\texttt{src/P01\_Proteccion\_Laser.py} (§\ref{sec:software}) permiten evaluar
automáticamente las 12 lentes contra cualquier escenario láser.

\subsection{Valores de MPE de Referencia}

\begin{itemize}
	\item \textbf{L1} (532~nm, 8~ns, pulsado):
	      $H_{\mathrm{MPE}} = \SI{5.0e-7}{\joule\per\centi\metre\squared}$
	      (pulso único).
	\item \textbf{L2} (514,5~nm, CW, 0,25~s):
	      $H_{\mathrm{MPE}} = \SI{2.5e-3}{\watt\per\centi\metre\squared}$.
	      Exposición prolongada ($>$10~s):
	      $H_{\mathrm{MPE}} = \SI{1.0e-3}{\watt\per\centi\metre\squared}$.
	\item \textbf{L3} (635~nm, CW, 0,25~s):
	      $H_{\mathrm{MPE}} = \SI{2.5e-3}{\watt\per\centi\metre\squared}$.
\end{itemize}

% =============================================================================
\section{Metodología de Cálculo}
\label{sec:metodologia}

\subsection{Paso 1: Área del Haz}
\begin{equation}
	A = \frac{\pi\, a^2}{4} \quad\text{($a$ en cm $\Rightarrow$ $A$ en cm$^2$)}
	\label{eq:area}
\end{equation}

\subsection{Paso 2: Exposición Potencial}
\begin{itemize}
	\item CW: $H_0 = P/A$ \;[\si{\watt\per\centi\metre\squared}]
	\item Pulsado: $H_0 = E/A$ \;[\si{\joule\per\centi\metre\squared}]
\end{itemize}

\subsection{Paso 3: OD Requerida}
\begin{equation}
	OD_{\mathrm{req}} = \log_{10}\!\left(\frac{H_0}{H_{\mathrm{MPE}}}\right)
	\label{eq:OD_req}
\end{equation}

\subsection{Paso 4: Selección de EPO}
OD de la gafa $\ge OD_{\mathrm{req}}$ a la $\lambda$ del láser. Si varias
cumplen, preferir mayor VLT.

% =============================================================================
\section{Ejemplo Resuelto: Tres Escenarios Base}
\label{sec:calculos}

\subsection{L1: PIV -- Nd:YAG 532~nm (Clase~4)}

$a = \SI{5}{\mm} = \SI{0.5}{\cm}$ $\Rightarrow$
$A = \pi(0{,}5)^2/4 = \SI{0.1963}{\centi\metre\squared}$

$E = \SI{200}{\mJ} = \SI{0.2}{\joule}$:
\begin{align}
	H_0               & = \frac{0{,}2}{0{,}1963} \approx \SI{1.02}{\joule\per\centi\metre\squared}
	\label{eq:H0_L1}                                                                               \\
	OD_{\mathrm{req}} & = \log_{10}\!\left(\frac{1{,}02}{5\times10^{-7}}\right)
	\approx \mathbf{6{,}31}
	\label{eq:OD_L1}
\end{align}

\paragraph{Evaluación de las 12 lentes a 532~nm (Tabla~\ref{tab:eval_L1}).}

\begin{table}[H]
	\centering
	\caption{Evaluación del catálogo para L1 -- PIV 532~nm ($OD_{\mathrm{req}} = 6{,}31$).}
	\label{tab:eval_L1}
	\small
	\begin{tabular}{c c c c c l}
		\toprule
		\textbf{Lente} & \textbf{OD @ 532} & \textbf{$\ge OD_{\mathrm{req}}$?}
		               & \textbf{Margen}   & \textbf{VLT}                      & \textbf{Veredicto}                                                               \\
		\midrule
		L-01           & 5                 & No                                & $-1{,}31$          & 40\% & \textcolor{aerored}{Insuficiente}                    \\
		L-03           & 6                 & No                                & $-0{,}31$          & 15\% & \textcolor{aerored}{Insuficiente}                    \\
		\rowcolor{aerogreen!15}
		L-05           & 7                 & Sí                                & $+0{,}69$          & 30\% & \textcolor{aerogreen}{\textbf{Adecuada}}             \\
		L-06           & 3                 & No                                & $-3{,}31$          & 50\% & \textcolor{aerored}{Insuficiente}                    \\
		L-07           & 4                 & No                                & $-2{,}31$          & 35\% & \textcolor{aerored}{Insuficiente}                    \\
		L-09           & 2                 & No                                & $-4{,}31$          & 45\% & \textcolor{aerored}{Insuficiente}                    \\
		\rowcolor{aerogreen!15}
		L-10           & 6                 & Sí*                               & $-0{,}31$          & 20\% & \textcolor{aeroyellow}{Marginal (6+ $\approx$ 6,31)} \\
		L-12           & 4                 & No                                & $-2{,}31$          & 22\% & \textcolor{aerored}{Insuficiente}                    \\
		\midrule
		\multicolumn{6}{l}{\footnotesize L-02, L-04, L-08, L-11: sin cobertura a 532~nm $\rightarrow$ \textcolor{aerored}{No aplica}.}                            \\
		\bottomrule
	\end{tabular}
\end{table}

\textbf{Selección: L-05} (OD~7, margen +0,69; VLT 30\%).
L-10 es marginal (OD~6+ podría no alcanzar 6,31).

\subsection{L2: LDA -- Ar-Ion 514,5~nm (Clase~4)}

$a = \SI{1.2}{\mm} = \SI{0.12}{\cm}$ $\Rightarrow$
$A = \pi(0{,}12)^2/4 = \SI{1.131e-2}{\centi\metre\squared}$

$P = \SI{1.5}{\watt}$:
\begin{align}
	H_0               & = \frac{1{,}5}{1{,}131\times10^{-2}} \approx \SI{132.6}{\watt\per\centi\metre\squared}
	\label{eq:H0_L2}                                                                                           \\
	OD_{\mathrm{req}} & = \log_{10}\!\left(\frac{132{,}6}{2{,}5\times10^{-3}}\right)
	\approx \mathbf{4{,}72}
	\label{eq:OD_L2}
\end{align}

\paragraph{Evaluación a 514,5~nm (Tabla~\ref{tab:eval_L2}).}

\begin{table}[H]
	\centering
	\caption{Evaluación del catálogo para L2 -- LDA 514,5~nm ($OD_{\mathrm{req}} = 4{,}72$).}
	\label{tab:eval_L2}
	\small
	\begin{tabular}{c c c c c l}
		\toprule
		\textbf{Lente} & \textbf{OD @ 514,5} & \textbf{$\ge OD_{\mathrm{req}}$?}
		               & \textbf{Margen}     & \textbf{VLT}                      & \textbf{Veredicto}                                                     \\
		\midrule
		\rowcolor{aerogreen!15}
		L-01           & 5                   & Sí                                & $+0{,}28$          & 40\% & \textcolor{aerogreen}{\textbf{Adecuada}}   \\
		L-03           & 7                   & Sí                                & $+2{,}28$          & 15\% & \textcolor{aerogreen}{Adecuada (baja VLT)} \\
		\rowcolor{aerogreen!15}
		L-05           & 7                   & Sí                                & $+2{,}28$          & 30\% & \textcolor{aerogreen}{\textbf{Adecuada}}   \\
		L-06           & 3                   & No                                & $-1{,}72$          & 50\% & \textcolor{aerored}{Insuficiente}          \\
		L-07           & 4                   & No                                & $-0{,}72$          & 35\% & \textcolor{aerored}{Insuficiente}          \\
		L-09           & 2                   & No                                & $-2{,}72$          & 45\% & \textcolor{aerored}{Insuficiente}          \\
		\rowcolor{aerogreen!15}
		L-10           & 6                   & Sí                                & $+1{,}28$          & 20\% & \textcolor{aerogreen}{Adecuada}            \\
		L-12           & 4                   & No                                & $-0{,}72$          & 22\% & \textcolor{aerored}{Insuficiente}          \\
		\midrule
		\multicolumn{6}{l}{\footnotesize L-02, L-04, L-08, L-11: sin cobertura a 514,5~nm $\rightarrow$ \textcolor{aerored}{No aplica}.}                  \\
		\bottomrule
	\end{tabular}
\end{table}

\textbf{Selección: L-01} (OD~5, margen +0,28; mejor VLT = 40\%).

\subsection{L3: Alineación -- Diodo 635~nm (Clase~3R)}

$a = \SI{3}{\mm} = \SI{0.3}{\cm}$ $\Rightarrow$
$A = \pi(0{,}3)^2/4 = \SI{7.069e-2}{\centi\metre\squared}$

$P = \SI{4.5}{\mW} = \SI{4.5e-3}{\watt}$:
\begin{align}
	H_0               & = \frac{4{,}5\times10^{-3}}{7{,}069\times10^{-2}}
	\approx \SI{6.37e-2}{\watt\per\centi\metre\squared}
	\label{eq:H0_L3}                                                                             \\
	OD_{\mathrm{req}} & = \log_{10}\!\left(\frac{6{,}37\times10^{-2}}{2{,}5\times10^{-3}}\right)
	\approx \mathbf{1{,}41}
	\label{eq:OD_L3}
\end{align}

\paragraph{Evaluación a 635~nm (Tabla~\ref{tab:eval_L3}).}

\begin{table}[H]
	\centering
	\caption{Evaluación del catálogo para L3 -- Alineación 635~nm ($OD_{\mathrm{req}} = 1{,}41$).}
	\label{tab:eval_L3}
	\small
	\begin{tabular}{c c c c c l}
		\toprule
		\textbf{Lente} & \textbf{OD @ 635} & \textbf{$\ge OD_{\mathrm{req}}$?}
		               & \textbf{Margen}   & \textbf{VLT}                      & \textbf{Veredicto}                                                     \\
		\midrule
		\rowcolor{aerogreen!15}
		L-02           & 2                 & Sí                                & $+0{,}59$          & 55\% & \textcolor{aerogreen}{\textbf{Adecuada}}   \\
		\rowcolor{aerogreen!15}
		L-08           & 3                 & Sí                                & $+1{,}59$          & 25\% & \textcolor{aerogreen}{Adecuada}            \\
		L-11           & 1                 & No                                & $-0{,}41$          & 60\% & \textcolor{aerored}{Insuficiente}          \\
		\rowcolor{aerogreen!15}
		L-12           & 5                 & Sí                                & $+3{,}59$          & 22\% & \textcolor{aerogreen}{Adecuada (excesiva)} \\
		\midrule
		\multicolumn{6}{l}{\footnotesize L-01, L-03, L-04, L-05, L-06, L-07, L-09, L-10: sin cobertura a 635~nm $\rightarrow$ N/A.}                     \\
		\bottomrule
	\end{tabular}
\end{table}

\textbf{Selección: L-02} (OD~2, margen +0,59; mejor VLT = 55\%).

\subsection{Tabla Resumen}

\begin{table}[H]
	\centering
	\caption{Resumen de resultados para los tres escenarios base.}
	\label{tab:resumen}
	\small
	\begin{tabular}{l c c c c c}
		\toprule
		\textbf{Láser (Clase)}    & $\boldsymbol{\lambda}$
		                          & $\boldsymbol{H_0}$               & $\boldsymbol{H_{\mathrm{MPE}}}$
		                          & $\boldsymbol{OD_{\mathrm{req}}}$ & \textbf{Lente}                                                                    \\
		\midrule
		L1 -- PIV (\textbf{4})    & 532~nm                           & 1,02~J/cm$^2$                   & $5\!\times\!10^{-7}$~J/cm$^2$     & 6,31 & L-05 \\
		L2 -- LDA (\textbf{4})    & 514,5~nm                         & 132,6~W/cm$^2$                  & $2{,}5\!\times\!10^{-3}$~W/cm$^2$ & 4,72 & L-01 \\
		L3 -- Alin. (\textbf{3R}) & 635~nm                           & 0,064~W/cm$^2$                  & $2{,}5\!\times\!10^{-3}$~W/cm$^2$ & 1,41 & L-02 \\
		\bottomrule
	\end{tabular}
\end{table}

\subsection{Evaluación de Riesgo de los Escenarios Base}
\label{sec:riesgo_base}

Aplicando la matriz de riesgos 5$\times$5 (§\ref{sec:riesgos}), cada escenario
se evalúa según su probabilidad de exposición y severidad del daño potencial.

\begin{table}[H]
	\centering
	\caption{Evaluación de riesgo para los tres escenarios base.}
	\label{tab:riesgo_base}
	\small
	\begin{tabularx}{\textwidth}{l c c c c X}
		\toprule
		\textbf{Escenario} & \textbf{Prob.}                                   & \textbf{Sev.}                 & \textbf{$R$}
		                   & \textbf{Nivel}                                   & \textbf{Controles requeridos}                                                         \\
		\midrule
		\rowcolor{aerored!15}
		L1 -- PIV (Clase~4)
		                   & Alto (4)                                         & Mayor (4)                     & 16           & \textcolor{aerored}{\textbf{Alto}}
		                   & EPO obligatorio, interlock, señalización Clase~4                                                                                         \\
		\rowcolor{aerored!10}
		L2 -- LDA (Clase~4)
		                   & Medio (3)                                        & Mayor (4)                     & 12           & \textcolor{aeroyellow}{\textbf{Medio}}
		                   & EPO obligatorio, cortinas, SOPs                                                                                                          \\
		\rowcolor{aerogreen!10}
		L3 -- Alin. (Clase~3R)
		                   & MuyBajo (1)                                      & Menor (2)                     & 2            & \textcolor{aerogreen}{\textbf{Bajo}}
		                   & EPO recomendado, señalización                                                                                                            \\
		\bottomrule
	\end{tabularx}
\end{table}

\textbf{Justificación:}
\begin{itemize}
	\item \textbf{PIV:} potencia pulsada elevada (200~mJ), rango visible con
	      enfoque retiniano, operación en laboratorio abierto $\rightarrow$
	      probabilidad alta.
	\item \textbf{LDA:} haz CW colimado (1,5~W), visible, pero contenido en
	      banco óptico $\rightarrow$ probabilidad media.
	\item \textbf{Alineación:} baja potencia (4,5~mW), diámetro grande, Clase~3R
	      $\rightarrow$ riesgo inherente bajo.
\end{itemize}

% =============================================================================
\section{Casos de Estudio por Equipos}
\label{sec:equipos}

Cada equipo de trabajo recibe un escenario diferente que debe resolver de forma
autónoma. Los cuatro casos se basan en láseres reales empleados en laboratorios
aeronáuticos. Cada equipo debe:
\begin{enumerate}
	\item Calcular $A$, $H_0$ y $OD_{\mathrm{req}}$.
	\item Evaluar las 12 lentes del catálogo (Tabla~\ref{tab:epo}).
	\item Completar la tabla de evaluación (Tabla~\ref{tab:eval_template}).
	\item Seleccionar la lente óptima (OD suficiente + mejor VLT), o bien
	      justificar por qué ninguna cumple.
	\item Realizar la evaluación de riesgo: completar la tabla de riesgo
	      (Tabla~\ref{tab:risk_template}) asignando probabilidad (1--5),
	      severidad (1--5), calculando $R = P \times S$ y determinando el
	      nivel de riesgo según la matriz (§\ref{sec:riesgos}).
	\item Verificar con el script Python o el notebook.
	\item Resolver los tres problemas adicionales asignados.
\end{enumerate}

\begin{table}[H]
	\centering
	\caption{Plantilla de evaluación de lentes (a completar por cada equipo).}
	\label{tab:eval_template}
	\small
	\begin{tabular}{c c c c c l}
		\toprule
		\textbf{Lente} & \textbf{OD @ $\lambda$} & \textbf{$\ge OD_{\mathrm{req}}$?}
		               & \textbf{Margen}         & \textbf{VLT}                      & \textbf{Veredicto}      \\
		\midrule
		L-01           &                         &                                   &                    &  & \\
		L-02           &                         &                                   &                    &  & \\
		L-03           &                         &                                   &                    &  & \\
		L-04           &                         &                                   &                    &  & \\
		L-05           &                         &                                   &                    &  & \\
		L-06           &                         &                                   &                    &  & \\
		L-07           &                         &                                   &                    &  & \\
		L-08           &                         &                                   &                    &  & \\
		L-09           &                         &                                   &                    &  & \\
		L-10           &                         &                                   &                    &  & \\
		L-11           &                         &                                   &                    &  & \\
		L-12           &                         &                                   &                    &  & \\
		\bottomrule
	\end{tabular}
\end{table}

\begin{table}[H]
	\centering
	\caption{Plantilla de evaluación de riesgo (a completar por cada equipo).}
	\label{tab:risk_template}
	\small
	\begin{tabularx}{\textwidth}{l c X}
		\toprule
		\textbf{Factor}                   & \textbf{Valor (1--5)} & \textbf{Justificación}                                           \\
		\midrule
		Probabilidad ($P$)                &                       & (MuyBajo=1, Bajo=2, Medio=3, Alto=4, MuyAlto=5)                  \\
		Severidad ($S$)                   &                       & (Insignificante=1, Menor=2, Moderado=3, Mayor=4, Catastrófico=5) \\
		\midrule
		Riesgo $R = P \times S$           &                       &                                                                  \\
		Nivel de riesgo                   &                       & Bajo ($<5$) / Medio (5--12) / Alto ($>12$)                       \\
		\midrule
		Controles requeridos              & {--}                  & (Listar: EPO, ingeniería, administrativos, otros)                \\
		Posición en la matriz (fila, col) & {--}                  & (Marcar en Figura~\ref{fig:risk_matrix})                         \\
		\bottomrule
	\end{tabularx}
\end{table}

% ---- EQUIPO 1 ----
\subsection{Equipo~1: Nd:YAG Fundamental 1064~nm (Clase~4, Pulsado)}
\label{sec:eq1}

\begin{table}[H]
	\centering
	\caption{Parámetros del láser -- Equipo~1.}
	\small
	\begin{tabular}{l l}
		\toprule
		\textbf{Parámetro} & \textbf{Valor}                                                 \\
		\midrule
		Tipo / Aplicación  & Nd:YAG fundamental / PIV-IR, LIBS                              \\
		$\lambda$          & 1064~nm                                                        \\
		Modo               & Pulsado                                                        \\
		Energía $E$        & 500~mJ/pulso                                                   \\
		Duración $\tau$    & 10~ns                                                          \\
		Repetición $f$     & 10~Hz                                                          \\
		Diámetro $a$       & 7~mm                                                           \\
		\rowcolor{aerored!15}
		Clase              & \textbf{4}                                                     \\
		$H_{\mathrm{MPE}}$ & $\SI{5.0e-6}{\joule\per\centi\metre\squared}$ (pulso, 1064~nm) \\
		\bottomrule
	\end{tabular}
\end{table}

\textbf{Tareas:} Calcule $A$, $H_0$ y $OD_{\mathrm{req}}$. Complete la tabla
de evaluación para las 12 lentes a 1064~nm. Determine si alguna lente cumple o
si debe activarse el protocolo de carencia de EPO.

\textbf{Evaluación de riesgo:} complete la Tabla~\ref{tab:risk_template} para
este escenario. Considere que el haz es infrarrojo (invisible), pulsado con
alta energía y Clase~4. Ubique el escenario en la matriz de riesgos
(Figura~\ref{fig:risk_matrix}) y justifique los controles necesarios.

\subsubsection*{Problemas Adicionales -- Equipo~1}

\begin{enumerate}
	\item[\textbf{P1.1}] Si la energía del láser se reduce a 100~mJ/pulso
	      manteniendo el resto de parámetros, ¿cuánto cambia $OD_{\mathrm{req}}$?
	      ¿Alguna lente del catálogo cumple ahora?

	\item[\textbf{P1.2}] La reflexión especular del haz de 1064~nm sobre una
	      superficie metálica conserva el 60\% de la energía original. Calcule
	      $OD_{\mathrm{req}}$ para la reflexión y determine si las gafas
	      seleccionadas (o la falta de ellas) siguen siendo válidas.

	\item[\textbf{P1.3}] El mismo Nd:YAG genera simultáneamente el segundo
	      armónico a 532~nm con $E = 50$~mJ y $a = 4$~mm. Calcule
	      $OD_{\mathrm{req}}$ a 532~nm (use la MPE de L1 del ejemplo resuelto)
	      y proponga la combinación de protección necesaria para trabajar con
	      ambas longitudes de onda simultáneamente.
\end{enumerate}

% ---- EQUIPO 2 ----
\subsection{Equipo~2: CO$_2$ CW a 10\,600~nm (Clase~4)}
\label{sec:eq2}

\begin{table}[H]
	\centering
	\caption{Parámetros del láser -- Equipo~2.}
	\small
	\begin{tabular}{l l}
		\toprule
		\textbf{Parámetro} & \textbf{Valor}                                                      \\
		\midrule
		Tipo / Aplicación  & CO$_2$ / Corte industrial, DIAL                                     \\
		$\lambda$          & 10\,600~nm                                                          \\
		Modo               & CW                                                                  \\
		Potencia $P$       & 50~W                                                                \\
		Diámetro $a$       & 8~mm                                                                \\
		\rowcolor{aerored!15}
		Clase              & \textbf{4}                                                          \\
		$H_{\mathrm{MPE}}$ & $\SI{0.1}{\watt\per\centi\metre\squared}$ (CW, 10\,600~nm, $>10$~s) \\
		\bottomrule
	\end{tabular}
\end{table}

\textbf{Tareas:} Calcule $A$, $H_0$ y $OD_{\mathrm{req}}$. Evalúe las 12
lentes a 10\,600~nm. Indique cuántas ofrecen cobertura a esa longitud de onda
y si alguna cumple. Justifique el protocolo adecuado si ninguna es suficiente.

\textbf{Evaluación de riesgo:} complete la Tabla~\ref{tab:risk_template}.
Considere que el CO$_2$ tiene muy alta potencia (50~W), el haz es IR lejano
(invisible, daño corneal no retiniano), y se usa en modo CW continuo.
Compare la severidad corneal vs.\ retiniana.

\subsubsection*{Problemas Adicionales -- Equipo~2}

\begin{enumerate}
	\item[\textbf{P2.1}] A 10\,600~nm el daño principal es corneal, no
	      retiniano. Explique por qué la MPE es relativamente alta
	      (\SI{0.1}{\watt\per\centi\metre\squared}) comparada con las MPE en
	      el rango visible (400--700~nm). ¿Significa que el láser CO$_2$ es
	      menos peligroso que un láser visible de igual potencia?

	\item[\textbf{P2.2}] Si la potencia del láser se reduce a 5~W (modo
	      alineación), recalcule $OD_{\mathrm{req}}$. ¿La lente L-04
	      ($OD = 2$ a 10\,600~nm) es ahora suficiente?

	\item[\textbf{P2.3}] El haz CO$_2$ se enfoca con una lente convergente a
	      un punto de $a = 0{,}5$~mm para corte. Con los 50~W originales,
	      calcule la nueva $H_0$ y $OD_{\mathrm{req}}$. Compare con el caso
	      sin enfocar y discuta la importancia del diámetro del haz.
\end{enumerate}

% ---- EQUIPO 3 ----
\subsection{Equipo~3: Diodo 808~nm CW (Clase~3B)}
\label{sec:eq3}

\begin{table}[H]
	\centering
	\caption{Parámetros del láser -- Equipo~3.}
	\small
	\begin{tabular}{l l}
		\toprule
		\textbf{Parámetro} & \textbf{Valor}                                                    \\
		\midrule
		Tipo / Aplicación  & Diodo semiconductor / Bombeo, iluminación IR                      \\
		$\lambda$          & 808~nm                                                            \\
		Modo               & CW                                                                \\
		Potencia $P$       & 300~mW                                                            \\
		Diámetro $a$       & 2~mm                                                              \\
		\rowcolor{aeroyellow!30}
		Clase              & \textbf{3B}                                                       \\
		$H_{\mathrm{MPE}}$ & $\SI{2.5e-3}{\watt\per\centi\metre\squared}$ (CW, 808~nm, 0,25~s) \\
		\bottomrule
	\end{tabular}
\end{table}

\textbf{Tareas:} Calcule $A$, $H_0$ y $OD_{\mathrm{req}}$. Complete la tabla
de evaluación a 808~nm. Identifique todas las lentes que cumplen y seleccione
la óptima justificando con el criterio OD + VLT.

\textbf{Evaluación de riesgo:} complete la Tabla~\ref{tab:risk_template}.
El diodo 808~nm es infrarrojo cercano (invisible), Clase~3B con potencia
moderada. Considere que el haz alcanza la retina (400--1400~nm) y evalúe la
probabilidad en un entorno de laboratorio con banco óptico cerrado.

\subsubsection*{Problemas Adicionales -- Equipo~3}

\begin{enumerate}
	\item[\textbf{P3.1}] El diodo se emplea para bombear un cristal Nd:YVO$_4$
	      que emite a 1064~nm con $P = 200$~mW y $a = 0{,}8$~mm. Calcule
	      $OD_{\mathrm{req}}$ para el haz de 1064~nm
	      ($H_{\mathrm{MPE}} = \SI{5.0e-3}{\watt\per\centi\metre\squared}$,
	      CW, 0,25~s). ¿La misma lente seleccionada para 808~nm protege
	      también a 1064~nm?

	\item[\textbf{P3.2}] Compare la VLT de todas las lentes que cumplen a
	      808~nm. Si el operador necesita distinguir colores en una pantalla
	      durante el experimento, ¿cuál es la VLT mínima recomendable?
	      Justifique la selección final considerando ergonomía.

	\item[\textbf{P3.3}] Si la exposición se prolonga a $>$10~s, la MPE a
	      808~nm se reduce a
	      $H_{\mathrm{MPE}} = \SI{1.0e-3}{\watt\per\centi\metre\squared}$.
	      Recalcule $OD_{\mathrm{req}}$ y determine si alguna de las lentes
	      que antes cumplían ahora se vuelve marginal o insuficiente.
\end{enumerate}

% ---- EQUIPO 4 ----
\subsection{Equipo~4: He-Ne CW a 632,8~nm (Clase~3R) + LDA Prolongada}
\label{sec:eq4}

Este equipo tiene un caso compuesto: primero evalúa un He-Ne de alineación y
luego reevalúa el L2 (LDA 514,5~nm) bajo exposición prolongada ($>$10~s,
$H_{\mathrm{MPE}} = \SI{1.0e-3}{\watt\per\centi\metre\squared}$).

\paragraph{Caso~A -- He-Ne 632,8~nm.}

\begin{table}[H]
	\centering
	\caption{Parámetros del láser -- Equipo~4 (Caso~A).}
	\small
	\begin{tabular}{l l}
		\toprule
		\textbf{Parámetro} & \textbf{Valor}                                        \\
		\midrule
		Tipo / Aplicación  & He-Ne / Alineación óptica                             \\
		$\lambda$          & 632,8~nm                                              \\
		Modo               & CW                                                    \\
		Potencia $P$       & 10~mW                                                 \\
		Diámetro $a$       & 1,5~mm                                                \\
		\rowcolor{aeroyellow!30}
		Clase              & \textbf{3R}                                           \\
		$H_{\mathrm{MPE}}$ & $\SI{2.5e-3}{\watt\per\centi\metre\squared}$ (0,25~s) \\
		\bottomrule
	\end{tabular}
\end{table}

\textbf{Tareas Caso~A:} Calcule $A$, $H_0$ y $OD_{\mathrm{req}}$ a 632,8~nm.
Complete la tabla de evaluación. Seleccione la lente óptima.

\textbf{Evaluación de riesgo (ambos casos):} complete una
Tabla~\ref{tab:risk_template} para el Caso~A y otra para el Caso~B. El He-Ne
es visible y de baja potencia (Clase~3R), mientras que la LDA prolongada es
Clase~4 con exposición extendida. Compare cómo la duración de exposición
afecta la probabilidad y, por tanto, el nivel de riesgo.

\paragraph{Caso~B -- LDA 514,5~nm con exposición prolongada.}

Reutilice los datos completos del láser L2 (Tabla~\ref{tab:lasers}), pero
considere exposición prolongada $>$10~s con
$H_{\mathrm{MPE}} = \SI{1.0e-3}{\watt\per\centi\metre\squared}$.

\textbf{Tareas Caso~B:} Calcule $OD_{\mathrm{req}}$ bajo la nueva MPE.
Compare con el $OD_{\mathrm{req}} = 4{,}72$ del ejemplo resuelto (0,25~s).
¿Alguna lente que cumplía antes ahora es insuficiente o marginal?

\subsubsection*{Problemas Adicionales -- Equipo~4}

\begin{enumerate}
	\item[\textbf{P4.1}] Para el Caso~A, si la potencia del He-Ne se duplica
	      a 20~mW (un modelo de mayor potencia), recalcule $OD_{\mathrm{req}}$.
	      ¿La lente seleccionada sigue siendo adecuada o debe cambiar?

	\item[\textbf{P4.2}] Si un operador necesita gafas que protejan
	      simultáneamente contra el He-Ne (632,8~nm) y el LDA (514,5~nm,
	      exposición prolongada), ¿existe alguna lente en el catálogo que
	      cumpla ambos requisitos de OD al mismo tiempo? Identifique las
	      candidatas y seleccione la óptima.

	\item[\textbf{P4.3}] El LSO propone instalar un obturador automático que
	      limita la exposición máxima a 0,25~s para el LDA. ¿Cómo afecta esto
	      la selección de EPO? Calcule la nueva $OD_{\mathrm{req}}$ y compare
	      el número de lentes adecuadas con y sin obturador. Discuta si el
	      control de ingeniería puede ampliar las opciones de protección.
\end{enumerate}

% =============================================================================
\section{Evaluación de Riesgos}
\label{sec:riesgos}

\subsection{Matriz de Riesgos}

\begin{figure}[H]
	\centering
	\begin{tikzpicture}[scale=0.85]
		\foreach \x in {0,...,4} {
				\foreach \y in {0,...,4} {
						\pgfmathsetmacro{\risk}{(\x+1)*(\y+1)}
						\pgfmathparse{\risk < 5 ? "aerogreen!40" : (\risk < 13 ? "aeroyellow!60" : "aerored!50")}
						\edef\fillcol{\pgfmathresult}
						\fill[\fillcol] (\x,\y) rectangle (\x+1,\y+1);
						\draw[gray] (\x,\y) rectangle (\x+1,\y+1);
						\node at (\x+0.5,\y+0.5) {\pgfmathprintnumber[fixed,precision=0]{\risk}};
					}
			}
		\foreach \x/\l in {0/MuyBajo,1/Bajo,2/Medio,3/Alto,4/MuyAlto} {
				\node[font=\tiny,rotate=45,anchor=west] at (\x+0.5,-0.1) {\l};
			}
		\foreach \y/\l in {0/Insignificante,1/Menor,2/Moderado,3/Mayor,4/Catastrófico} {
				\node[font=\tiny,anchor=east] at (-0.1,\y+0.5) {\l};
			}
		\node[font=\small\bfseries,rotate=90] at (-1.5,2.5) {Severidad};
		\node[font=\small\bfseries] at (2.5,-1.2) {Probabilidad};
		\fill[blue,opacity=0.8] (3.5,3.5) circle (0.25);
		\node[white,font=\tiny\bfseries] at (3.5,3.5) {PIV};
		\fill[blue!70,opacity=0.8] (2.5,3.5) circle (0.25);
		\node[white,font=\tiny\bfseries] at (2.5,3.5) {LDA};
		\fill[orange,opacity=0.8] (0.5,1.5) circle (0.25);
		\node[white,font=\tiny\bfseries] at (0.5,1.5) {Alin.};
		\fill[aerogreen!40] (5.5,3.5) rectangle (6.0,4.0);
		\node[right,font=\tiny] at (6.0,3.75) {Bajo ($<5$)};
		\fill[aeroyellow!60] (5.5,2.8) rectangle (6.0,3.3);
		\node[right,font=\tiny] at (6.0,3.05) {Medio (5--12)};
		\fill[aerored!50] (5.5,2.1) rectangle (6.0,2.6);
		\node[right,font=\tiny] at (6.0,2.35) {Alto ($>12$)};
	\end{tikzpicture}
	\caption{Matriz de riesgos. PIV y LDA (Clase~4): zona roja; Alineación
		(Clase~3R): zona verde.}
	\label{fig:risk_matrix}
\end{figure}

\subsection{Jerarquía de Controles}

\begin{enumerate}
	\item \textbf{Eliminación/Sustitución}: clase inferior si es posible.
	\item \textbf{Controles de ingeniería}: enclave, obturador, cabina, cortinas.
	\item \textbf{Controles administrativos}: señalización, LSO, SOPs, registro.
	\item \textbf{EPP}: gafas EN~207~\cite{en207} con OD y VLT adecuados.
\end{enumerate}

% =============================================================================
\section{Herramientas de Software}
\label{sec:software}

Carpeta \texttt{Practicas/P01\_Proteccion\_Laser/}:

\subsection{Notebook Jupyter}
\texttt{notebooks/P01\_Proteccion\_Laser.ipynb} -- cálculo interactivo de $H_0$
y OD para los 3 escenarios base, evaluación de riesgo con matriz 5$\times$5,
verificación contra catálogo CSV, gráficas de margen OD.

\subsection{Script Python}
\texttt{src/P01\_Proteccion\_Laser.py} -- funciones de cálculo reutilizables:

\begin{lstlisting}[language=Python, caption={Funciones principales.}]
import math
def area_cm2_from_diameter_mm(a_mm):
    d_cm = a_mm / 10.0
    return math.pi * (d_cm / 2.0) ** 2

def od_required(H0, HMPE):
    if H0 <= HMPE: return 0.0
    return math.log10(H0 / HMPE)

def risk_assessment(prob, severity):
    """Devuelve (R, nivel, color)."""
    R = prob * severity
    if R < 5:   return R, 'Bajo', 'green'
    if R < 13:  return R, 'Medio', 'yellow'
    return R, 'Alto', 'red'
\end{lstlisting}

\subsection{Datos}
\texttt{data/epo\_lenses.csv} -- catálogo de 12 lentes con columnas:
\texttt{lens\_id}, \texttt{brand}, \texttt{model}, \texttt{band\_lo\_nm},
\texttt{band\_hi\_nm}, \texttt{od\_value}, \texttt{vlt\_pct}.

% =============================================================================
\section{Procedimiento Experimental}
\label{sec:procedimiento}

\begin{enumerate}
	\item Identificar el láser asignado al equipo.
	\item Calcular $A$, $H_0$ y $OD_{\mathrm{req}}$.
	\item Evaluar las 12 lentes del catálogo.
	\item Seleccionar la lente óptima y justificar.
	\item Completar la evaluación de riesgo (Tabla~\ref{tab:risk_template}):
	      asignar probabilidad y severidad, calcular $R$, determinar nivel
	      y ubicar el escenario en la matriz (Figura~\ref{fig:risk_matrix}).
	\item Verificar con el notebook / script Python (incluyendo evaluación
	      de riesgo automatizada).
	\item Registrar resultados en el informe.
\end{enumerate}

% =============================================================================
\section{Cuestionario}
\label{sec:cuestionario}

\begin{enumerate}
	\item ¿Qué significa «OD~5+ @ 190--540~nm»? ¿En cuánto reduce la intensidad
	      a 532~nm? ¿Y a 600~nm?
	\item ¿Por qué un Clase~4 es peligroso incluso con reflexiones difusas?
	\item Mencione al menos tres parámetros necesarios para calcular
	      $OD_{\mathrm{req}}$.
	\item Entre dos lentes de igual OD, una con VLT~45\% y otra 20\%: ¿cuál
	      elegir y por qué?
	\item ¿Gafas para 1064~nm protegen a 532~nm del mismo Nd:YAG? Justifique.
	\item ¿Qué protocolo se sigue cuando ninguna lente del catálogo cumple
	      el $OD_{\mathrm{req}}$?
\end{enumerate}

% =============================================================================
\section{Formato del Informe}
\label{sec:informe}

Formato AIAA con secciones: Resumen, Introducción, Fundamento Teórico,
Metodología, Resultados (escenarios base + caso del equipo), Discusión
(cuestionario), Conclusiones, Referencias, Apéndice~A (gráficas Python).

% =============================================================================
\section{Conclusiones}
\label{sec:conclusiones}

\begin{enumerate}
	\item El catálogo de 12 lentes permite evaluar escenarios desde UV hasta
	      IR lejano, incluyendo situaciones donde \textbf{ninguna lente cumple}.
	\item La selección de EPO debe considerar simultáneamente: cobertura
	      espectral a $\lambda$ del láser, OD $\ge OD_{\mathrm{req}}$ y la
	      mayor VLT posible para garantizar la ergonomía.
	\item La duración de la exposición puede modificar la MPE y, por tanto,
	      la $OD_{\mathrm{req}}$: un mismo escenario puede requerir gafas
	      distintas según el tiempo de operación.
	\item Cuando el catálogo disponible es insuficiente, el protocolo es
	      informar al LSO y suspender operaciones hasta disponer de EPO
	      adecuado.
	\item Las herramientas Python automatizan todo el flujo de evaluación
	      (OD, selección de lentes y evaluación de riesgo) y facilitan la
	      verificación de los cálculos manuales.
\end{enumerate}

% =============================================================================
\printbibliography[title={Referencias}]

% =============================================================================
\appendix
\section{Tabla de MPE de Referencia}
\label{sec:appendix_mpe}

\begin{table}[H]
	\centering
	\caption{Valores de MPE para exposición ocular -- resumen (IEC~60825-1).}
	\label{tab:mpe_ref}
	\small
	\begin{tabular}{l l l l}
		\toprule
		\textbf{Rango $\lambda$} & \textbf{Duración}                 & \textbf{MPE}         & \textbf{Unidades} \\
		\midrule
		400--700~nm              & $<\SI{1.8e-5}{\second}$ (pulsado) & $5{,}0\times10^{-7}$ & J/cm$^2$          \\
		400--700~nm              & 0,25~s (reflejo aversión)         & $2{,}5\times10^{-3}$ & W/cm$^2$          \\
		400--700~nm              & 10~s (prolongada)                 & $1{,}0\times10^{-3}$ & W/cm$^2$          \\
		700--1050~nm             & 0,25~s                            & $2{,}5\times10^{-3}$ & W/cm$^2$          \\
		1050--1400~nm            & $<\SI{1e-5}{\second}$ (pulsado)   & $5{,}0\times10^{-6}$ & J/cm$^2$          \\
		$>$1400~nm               & $>$10~s (CW)                      & $1{,}0\times10^{-1}$ & W/cm$^2$          \\
		\bottomrule
	\end{tabular}
\end{table}

\end{document}
